\chapter{Kết luận và hướng phát triển}
\chaptergoal{Tổng hợp đóng góp học thuật và kỹ thuật của đề tài, đồng thời đề xuất lộ trình mở rộng có cơ sở cho các giai đoạn tiếp theo.}

\section{Kết luận chính của đề tài}
Đề tài đã thiết lập được một khung triển khai tương đối hoàn chỉnh cho recommendation dựa trên văn bản trong miền phim:
\begin{itemize}
    \item dữ liệu được chuẩn hóa thành lớp \texttt{core\_*} phục vụ trực tiếp cho NLP;
    \item pipeline huấn luyện được tách riêng trong thư mục \texttt{training/};
    \item kết quả mô hình được xuất thành artifact và nạp vào runtime;
    \item giao diện web đã có luồng discovery, chi tiết phim và recommendation.
\end{itemize}

Điểm quan trọng nhất là dự án đã chuyển từ cách làm ``chạy được'' sang cách làm ``có thể tái lập và mở rộng'', tức phù hợp với tinh thần nghiên cứu thực nghiệm.

\section{Giá trị phương pháp}
Từ góc nhìn phương pháp luận, đề tài cho thấy ba nguyên tắc có hiệu quả rõ rệt:
\begin{enumerate}[label=\arabic*.]
    \item \textbf{Data contract trước, mô hình sau}: kiểm soát chất lượng dữ liệu giúp giảm nhiễu thực nghiệm.
    \item \textbf{So sánh công bằng}: giữ cùng đơn vị biểu diễn để so sánh giữa các họ mô hình.
    \item \textbf{Artifact-first serving}: tách huấn luyện khỏi runtime giúp hệ thống ổn định hơn khi triển khai.
\end{enumerate}

\section{Những điểm còn hạn chế}
Mặc dù kiến trúc đã rõ ràng, chất lượng mô hình vẫn bị giới hạn bởi dữ liệu đầu vào:
\begin{itemize}
    \item nhiều phim chưa có review hoặc review quá ngắn;
    \item thiếu các trường metadata giàu thông tin cho xếp hạng;
    \item nhãn đánh giá retrieval mới ở mức weak supervision;
    \item chưa có dữ liệu tương tác người dùng để đánh giá cá nhân hóa.
\end{itemize}

\section{Hướng phát triển ngắn hạn}
Trong ngắn hạn, ưu tiên cao nhất là tăng độ tin cậy của kết quả thực nghiệm:
\begin{enumerate}[label=\arabic*.]
    \item hoàn thiện backfill dữ liệu metadata quan trọng;
    \item tăng chất lượng provenance và kiểm định ngôn ngữ review;
    \item chuẩn hóa bộ truy vấn đánh giá thủ công theo nhiều nhóm thể loại;
    \item hoàn tất bảng so sánh định lượng giữa baseline và mô hình cải tiến;
    \item cập nhật toàn bộ hình/bảng thực nghiệm vào báo cáo.
\end{enumerate}

\section{Hướng phát triển trung và dài hạn}
\begin{itemize}
    \item \textbf{Mở rộng biểu diễn}: bổ sung cast, crew, keywords, runtime vào movie profile.
    \item \textbf{Nâng cấp truy hồi}: chuyển từ full similarity matrix sang ANN indexing cho quy mô lớn.
    \item \textbf{Fine-tuning miền phim}: tinh chỉnh embedding theo tập review điện ảnh để tăng tính đặc thù miền.
    \item \textbf{Hybrid recommendation}: kết hợp content-based và tín hiệu hành vi người dùng.
    \item \textbf{Đánh giá online}: thu phản hồi click/like/watch để xây dựng vòng lặp cải thiện liên tục.
\end{itemize}

\section{Thông điệp tổng kết}
\ProjectName{} cho thấy một bài toán NLP ứng dụng có thể được phát triển theo hướng khoa học nếu mỗi quyết định kỹ thuật đều có bằng chứng dữ liệu đi kèm. Kết quả hiện tại chưa phải điểm kết thúc của hệ recommendation, nhưng đã hình thành một nền tảng đủ tốt để mở rộng thành một hệ thống gợi ý thực thụ với vòng đời nghiên cứu--triển khai rõ ràng.
